Сдвигает ли поле дислокацию, проверено прямым расчётом. Пластина алюминия
содержит то же закреплённое включение и прямую краевую дислокацию в плоскости скольжения
на 12~\AA{} выше включения (Дополнительные материалы, раздел~S6). Без
деформации включения дислокация после минимизации энергии стоит в «DISL_D»~\AA{} от
центра включения вдоль $x$, где среднее по линии касательное напряжение поля
около «DISL_TAU»~МПа. Когда включению задают деформацию~(\ref{eq:strain}),
дислокация отходит от него на «DISL_SHIFT»~\AA, примерно на половину
расстояния $b/2=1{,}43$~\AA{} между соседними долинами рельефа решётки
(положениями, где энергия дислокации наименьшая), и в соседнюю долину не
переходит. Когда деформацию снимают, дислокация возвращается на место
с точностью «DISL_REV»~\AA, то есть в этом расчёте при 0~К отклик упругий
и обратимый. Среднее по линии напряжение поля на дислокации примерно вдвое
меньше сопротивления решётки в используемом потенциале, около 9~МПа
(Дополнительные материалы, раздел~S6). Там, где сопротивление
меньше, например при комнатной температуре, то же поле могло бы сдвинуть
дислокацию дальше; в сплаве дислокации удерживают ещё атомы Mg и Si
и выделения, и сопротивление там выше (13--20~МПа).
